(a) The inverses $U_i\to f^{-1}(U_i)$ agree on overlaps, since each is an inverse to the same restricted morphism. They glue to an inverse of $f$.

(b) If $X$ is affine, take $f_1=1$. Conversely, write $1=\sum_i a_if_i$. No point can have all $(f_i)_x$ in its maximal ideal, so the $X_{f_i}$ cover $X$. Each intersection $X_{f_i}\cap X_{f_j}$ is a distinguished open in the affine scheme $X_{f_i}$ by \exref{II.2.16}(a), hence is affine and quasi-compact.

By \exref{II.2.4}, the identity of $A$ induces a morphism $f:X\to\operatorname{Spec}A$. Locality on stalks gives $f^{-1}(D(f_i))=X_{f_i}$. By \exref{II.2.16}(d), its restriction is the affine morphism corresponding to the canonical isomorphism $A_{f_i}\cong\Gamma(X_{f_i},\mathcal O_X)$, so it is an isomorphism onto $D(f_i)$. These opens cover $\operatorname{Spec}A$, and (a) applies.
